% Part: computability
% Chapter: computability-theory
% Section: universal-part-function

\documentclass[../../../include/open-logic-section]{subfiles}

\begin{document}

\olfileid{cmp}{thy}{uni}
\olsection{સર્વવ્યાપી આંશિક સંગણનીય વિધેય}

\begin{thm}
\ollabel{thm:univ-comp}
સર્વવ્યાપી આંશિક સંગણનીય વિધેય~$\fn{Un}(e,x)$ છે. બીજા શબ્દોમાં,
એવું વિધેય $\fn{Un}(e,x)$ છે કે:
\begin{enumerate}
\item $\fn{Un}(e,x)$ આંશિક સંગણનીય છે.
\item $f(x)$ કોઈ પણ આંશિક સંગણનીય વિધેય હોય, તો એવી
પ્રાકૃતિક સંખ્યા $e$ છે કે દરેક~$x$ માટે $f(x) \simeq \fn{Un}(e,x)$.
\end{enumerate}
\end{thm}

\begin{proof}
$\fn{Un}(e,x) \simeq U(\umin{s}{T(e,x,s)})$ રાખો, જ્યાં $U$ અને $T$
ક્લીનીના સામાન્ય સ્વરૂપના પ્રમેય (\olref[nfm]{thm:normal-form})માં
આપેલાં છે.
\end{proof}

\begin{explain}
આ તો આંશિક સંગણનીય વિધેયોની અસરકારક પરિગણના છે એવું ચોક્કસ રીતે
કહે છે. વિચાર એવો છે કે $f_e(x) = \fn{Un}(e,x)$ વડે વ્યાખ્યાયિત
વિધેયને $f_e$ લખીએ, તો અનુક્રમ $f_0$, $f_1$, $f_2$, \dots માં બધાં
આંશિક સંગણનીય વિધેયો આવે છે. વધુમાં $f_e(x)$ની ગણતરી $e$ અને~$x$
બંનેની દૃષ્ટિએ ``એકસરખી રીતે'' થઈ શકે છે. સરળતા ખાતર અહીં એક-આર્ગ્યુમેન્ટવાળાં
વિધેયો માટે સર્વવ્યાપી દ્વિ-આર્ગ્યુમેન્ટવાળું વિધેય લીધું છે; પરંતુ
સંખ્યાઓના અનુક્રમોને કોડ કરીને આ વાત વધુ આર્ગ્યુમેન્ટો માટે સહેલાઈથી
વિસ્તારી શકાય છે. ઉદાહરણ તરીકે, $f(x,y,z)$ જો $3$-સ્થાનીય આંશિક
પુનરાવર્તી વિધેય હોય, તો $g(x) \simeq f((x)_0, (x)_1, (x)_2)$
એક-આર્ગ્યુમેન્ટવાળું પુનરાવર્તી વિધેય છે.
\end{explain}

\end{document}
